% EXPECTED-LEAKS: affiliation, ead, cortext, cormark, fnmark, fntext, author
\documentclass[preprint,12pt]{elsarticle}

\journal{Journal of Computational Methods}

\begin{document}

\begin{frontmatter}

\title{Adaptive Mesh Refinement for Shock-Dominated Flows}

\author[lab1]{Hana Ito\corref{cor1}}
\ead{hana.ito@example.edu}
\author[lab2]{Luka Petrovic}
\cortext[cor1]{Corresponding author}
\affiliation[lab1]{organization={Institute of Applied Mathematics},addressline={City, Country}}
\affiliation[lab2]{organization={Research Centre}}

\begin{abstract}
We present an adaptive mesh refinement strategy that reduces compute requirements by 3x on shock-dominated flow problems while preserving accuracy.
\end{abstract}

\begin{keyword}
mesh refinement \sep shock capture \sep computational fluid dynamics
\end{keyword}

\end{frontmatter}

\section{Introduction}
Shock-dominated flows remain computationally demanding. Uniform meshes waste resolution in smooth regions.

\section{Method}
Our refinement indicator combines gradient magnitude with vorticity, weighted by local cell size.

\section{Results}
On the double Mach reflection benchmark, we achieve comparable solution quality to a uniform mesh with one third the cell count.

\section{Conclusion}
Adaptive refinement is a promising path for cost reduction in large-scale CFD simulations.

\end{document}
